\section{The diagonal Tornheim derivative in harmonic coordinates}
\label{sec:diagonal}

The incoming coincident report relates the digamma cube to a third
diagonal Tornheim derivative. We give the relevant derivation, with
its finite normalization visible, and then insert the new bridge.
This separates the earlier spectral identity from the additional
harmonic reduction proved here.

Write
\[
 \omega(t)=\T(t,t,t),\qquad \Omega_3=\omega'''(0),\qquad
 E(t)=\Gamma(1+t)e^{-Lt},\qquad
 h(t)=\zeta(1-t)+\frac1t.
\]
Here $\T$ is the continuation of
$\sum_{m,n\ge1}m^{-A}n^{-B}(m+n)^{-C}$, and the diagonal is
restricted before taking the limit at zero. The third derivative is
one-variable data; joint holomorphy at the origin is not assumed.

\subsection{The completed cubic generating function}

Set
\[
 g(t,x)=\zeta(1-t,x)+\frac1t
       =\sum_{m\ge0}\gamma_m(x)\frac{t^m}{m!}.
\]
Near $x=0$ its exact Hurwitz splitting yields
\begin{equation}\label{diag:local}
 g(t,x)=x^{t-1}+h(t)-(1-t)\zeta(2-t)x+O(x^2).
\end{equation}
After cubing, the only local monomials whose Mellin denominators
vanish at $t=0$ are
\[
 3h(t)^2x^{t-1},\qquad
 -3(1-t)\zeta(2-t)x^{2t-1}.
\]
Their complete spectral integrals are $3h(t)^2/t$ and
$-3(1-t)\zeta(2-t)/(2t)$. The complete numerators, rather than
just their values at zero, have to be removed to obtain the
unit-coordinate finite parts.

For $\Re t>1$, the absolutely convergent Hurwitz Fourier expansion
has coefficients
\[
 \widehat{\zeta(1-t,\cdot)}(n)
 =\Gamma(t)(2\pi|n|)^{-t}
             e^{-\pi i t\operatorname{sgn}(n)/2},\qquad n\ne0.
\]
The constant coefficient vanishes. In the product of three series,
the condition $n_1+n_2+n_3=0$ has six sign cones; on each, the
largest magnitude is the sum of the other two. Consequently
\begin{align}
 \int_0^1\zeta(1-t,x)^3\dd x
 &=6\Gamma(t)^3(2\pi)^{-3t}\cos(\pi t/2)\omega(t),\notag\\
 \int_0^1\zeta(1-t,x)^2\dd x
 &=2\Gamma(t)^2(2\pi)^{-2t}\zeta(2t).
 \label{diag:Fourier}
\end{align}
The classical Hurwitz--Tornheim connection is recorded, with its
literature attribution, in the coincident report
\cite{IncomingCoincident}. Finite endpoint Taylor subtraction
continues both sides meromorphically.

\begin{lemma}[Completed diagonal kernel]\label{lem:diag-completed}
The expression
\begin{align}
 F(t)={}&\frac{6E(t)^3\cos(\pi t/2)\omega(t)
                         +6E(t)^2\zeta(2t)+1}{t^3}\notag\\
 &-\frac{3h(t)^2}{t}
           +\frac{3(1-t)\zeta(2-t)}{2t}
 \label{diag:completed}
\end{align}
is holomorphic near zero and satisfies $F(0)=-Q$.
\end{lemma}

\begin{proof}
Expand $g^3$ and use \eqref{diag:Fourier} and the vanishing
one-factor integral. This gives the fraction on the first line as
the meromorphic integral of $g^3$. To justify its completion,
split the $x$ integral at any $b\in(0,1)$. Subtract finitely many
Taylor monomials of the smooth factor in \eqref{diag:local} on
$(0,b)$. The remainder and all fixed parameter derivatives are
integrable uniformly in a small $t$ disk. A local monomial with
exponent $kt-1$ has spectral integral $b^{kt}/(kt)$, while the
generating function of its coordinate finite parts is
$(b^{kt}-1)/(kt)$. Their difference is exactly $1/(kt)$.
All nonresonant denominators remain nonzero at zero. The two
resonant terms displayed above therefore give precisely the second
line of \eqref{diag:completed}. The resulting Taylor coefficients
are the finite parts of those of $g^3$. At zero,
$g(0,x)=\gamma_0(x)=-\psi(x)$, hence $F(0)=-Q$.
\end{proof}

Solving \eqref{diag:completed} for $\omega$ also proves its
holomorphy at zero, since the coefficient $E(t)^3\cos(\pi t/2)$
is nonzero there. This argument does not infer a joint regularity
statement from a ray restriction.

\subsection{An explicit reduction of the third diagonal derivative}

Use the classical values $\zeta(0)=-1/2$, $\zeta'(0)=-L/2$ and
\[
 E(t)=\exp\left[-(\gamma+L)t+\frac{\zeta(2)}2t^2
                              -\frac{\zeta(3)}3t^3+O(t^4)\right].
\]
The coefficients of $t^{-3},t^{-2},t^{-1}$ in
\eqref{diag:completed} give
\begin{equation}\label{diag:lower}
 \omega(0)=\frac13,\qquad \omega'(0)=L,
 \qquad \omega''(0)=L^2-4\zeta''(0).
\end{equation}
For clarity, after the first two substitutions the coefficient of
$t^2$ in the numerator is
\[
 3\omega''(0)+3\gamma^2-3L^2-\tfrac32\zeta(2)+12\zeta''(0),
\]
and the remaining two terms contribute
$-3\gamma^2+3\zeta(2)/2$ to the residue.
The numerator coefficient at $t^3$, after \eqref{diag:lower}, is
\[
 \Omega_3+8\zeta'''(0)+12(\gamma+L)\zeta''(0)
 +L^3+6\gamma L^2-5\gamma^3+\tfrac32\gamma\zeta(2).
\]
Since $h(t)=\gamma+\gamma_1t+O(t^2)$, the other constant
terms are $-6\gamma\gamma_1-3(\zeta(2)+\zeta'(2))/2$.
Thus the earlier cubic identity is recovered:
\begin{align}
 Q={}&-\Omega_3-8\zeta'''(0)-12(\gamma+L)\zeta''(0)
            -L^3-6\gamma L^2+5\gamma^3\notag\\
 &+6\gamma\gamma_1-\frac32\gamma\zeta(2)
                       +\frac32\bigl(\zeta(2)+\zeta'(2)\bigr).
 \label{diag:old-cube}
\end{align}

\begin{theorem}[Harmonic coordinate for the cubic diagonal]
\label{thm:diag-harmonic}
The third derivative of the specified diagonal continuation is
\begin{align}
 \Omega_3={}&6\eta(1)-8\zeta'''(0)
 -12(\gamma+L)\zeta''(0)-L^3-6\gamma L^2+5\gamma^3\notag\\
 &+12\gamma\gamma_1-\frac32\gamma\zeta(2)
                     -\frac32\zeta(2)+\frac32\zeta'(2).
 \label{bridge:tornheim}
\end{align}
In particular, \eqref{bridge:arithmetic} gives an absolutely
convergent generalized-harmonic evaluation of this derivative.
\end{theorem}
\begin{proof}
Substitute \eqref{bridge:eta} into \eqref{diag:old-cube} and solve
for $\Omega_3$.
\end{proof}

The independently certified arithmetic evaluation in the package gives
\begin{align*}
 Q&=1.96626273439153434066437533841163142382586897649\ldots,\\
 \eta(1)&=0.53678702435841856750360788434881167385105703307\ldots,\\
 \Omega_3&=86.6578666011000736529032768587621836100797302472\ldots.
\end{align*}
The first harmonic coefficient is therefore sufficient for the
diagonal cubic problem. A finite expression for that coefficient in
a specified algebra of ordinary zeta, Gamma and Stieltjes constants
remains an arithmetic research question.
